\section{Synthesis propositions and obstructions}\label{sec:register}
\cref{tab:register} lists the nine synthesis propositions with their status and what each rests on; the machine-readable register is \texttt{crosswalk/SYNTHESIS\_PROPOSITIONS.tsv}. A proposition is marked \emph{new} when its statement is derived here with an independent proof; no result is silently upgraded from a finite-exact status.
\begin{table}[t]\centering\footnotesize\setlength{\tabcolsep}{3pt}\begin{tabular}{@{}p{1.3cm}p{4.6cm}p{4.6cm}p{2.5cm}@{}}\toprule
ID & Subject & Status & Statement \\\midrule
SIV-P1 & Endpoint projection and retained information & NEW-PROVED & \cref{prop:P1}\\
SIV-P2 & Occurrence projection and primitive-relation information & NEW-PROVED (a); CITED (b),(c) & \cref{prop:P2}\\
SIV-P3 & Arithmetic availability versus geometric realization & NEW-PROVED (a),(b); CITED (c),(d) & \cref{prop:P3}\\
SIV-P4 & Arithmetic constraints and global history connectivity & NEW-PROVED from cited theorems; FINITE check & \cref{prop:P4}\\
SIV-P5 & Limits of endpoint geometry as a history invariant & NEW-PROVED (exhaustive recursion); FINITE-EXACT census & \cref{prop:P5}\\
SIV-P6 & Hierarchy of structural persistence under evolution & NEW-PROVED; (d) quotes finite-exact census & \cref{prop:P6}\\
SIV-P7 & Information requirements for topological prediction & NEW-PROVED (a)-(d); (e) search outcome only & \cref{prop:P7}\\
SIV-P8 & Scaling information versus geometric structure & NEW; FINITE-EXACT (k<=14, 12 periods) & \cref{prop:P8}\\
SIV-P9 & Correspondence levels for distinct arithmetic dynamics & NEW-PROVED (c); ASSEMBLED (a),(b),(d) & \cref{prop:P9}\\
\bottomrule\end{tabular}\caption{\textbf{Register of the synthesis propositions.} New results are proved in the text or in Appendix~\ref{app:proofs}; inherited ingredients are cited by theorem number (Appendix~\ref{app:prov}); finite-exact statements are labelled with their checked ranges.}\label{tab:register}\end{table}

\subsection{Answers to S1--S4}
\begin{description}[leftmargin=1.2em,itemsep=2pt]
\item[S1] \emph{Proved for integer states, open beyond.} The generative guard $(N+1)^2\mid t$ and the occurrence guard $d^2\mid t$ coincide under $d=N+1$ (\cref{prop:P2}(a)). For rational states the guard $d\equiv1\pmod b$ has no generative counterpart; a rational extension of the generative rule reproducing it is open (O8).
\item[S2] \emph{Proved, with obstructions.} The endpoint retains the valuation displacement $V_Ac$ and nothing finer (\cref{prop:P1}); the occurrence set retains everything that enters $\Kmove$ and nothing about the primitive system or about the number of histories (\cref{prop:P2,prop:P5}). Reconstruction of a history needs the fibre sequence along the path (\cref{thm:faith-cited}).
\item[S3] \emph{Proved obstruction and open maps.} $K(D(t))$ does not determine the topology of $\Hist(m_0,t)$ (\cref{prop:P5}); the route-complex side has independent arithmetic (\cref{prop:P4}); explicit maps between history geometry and evolving occurrence geometry beyond L2 are open (\cref{prop:P9}, O3).
\item[S4] \emph{Conditional (structural).} Both criteria are instances of monotone minimality (\cref{cor:S4}); an exact implication between them requires an identification of ambient orders and is open (O7).
\end{description}
